% Part: computability
% Chapter: computability-theory
% Section: universal-part-function

\documentclass[../../../include/open-logic-section]{subfiles}

\begin{document}

\olfileid[fr]{cmp}{thy}{uni}
\olsection{La fonction partielle calculable universelle}

\begin{thm}
\ollabel{thm:univ-comp}
Il existe une fonction partielle calculable
universelle~$\fn{Un}(e,x)$. Autrement dit,
il existe une fonction $\fn{Un}(e,x)$ telle que :
\begin{enumerate}
\item $\fn{Un}(e,x)$ est partielle calculable.
\item Pour toute fonction partielle calculable
  unaire~$f(x)$, il existe un entier naturel~$e$
  tel que $f(x) \simeq \fn{Un}(e,x)$ pour tout~$x$.
\end{enumerate}
\end{thm}

\begin{proof}
Posons $\fn{Un}(e,x) \simeq U(\umin{s}{T(e,x,s)})$,
où $U$ et $T$ sont ceux du théorème de forme
normale de Kleene
(\olref[nfm]{thm:normal-form}).
\end{proof}

\begin{explain}
Le théorème précise que les fonctions partielles
calculables admettent une énumération effective.
Si $f_e$ désigne la fonction définie par
$f_e(x) \simeq \fn{Un}(e,x)$, la suite
$f_0,f_1,f_2,\dots$ comprend toutes les fonctions
partielles calculables, et la valeur de $f_e(x)$
se calcule « uniformément » à partir de $e$
et de~$x$, lorsqu'elle est définie. Pour
simplifier, nous avons pris une fonction binaire
universelle pour les fonctions unaires; en codant
les suites de nombres, on peut étendre cette
construction à un nombre quelconque d'arguments.
Par exemple, si $f(x,y,z)$ est une fonction
récursive partielle à trois arguments, alors
$g(x) \simeq f((x)_0,(x)_1,(x)_2)$ est une
fonction récursive partielle unaire.\footnote{Correction éditoriale : la source appelle ici $g$ une fonction récursive unaire sans préciser sa partialité. Celle-ci doit être conservée après le codage des arguments. Nous employons aussi $\simeq$ dans la définition de $f_e$ pour couvrir les valeurs indéfinies.}
\end{explain}

\end{document}
